% Values: FAITHFUL, ILLUMINATING
% Evidence planned:
% - FAITHFUL: each row states what this definition decrees, with the
%   direction or ruling behind it (D2, D3, D5, D6, D7, D10, D11, A1, A11, A12,
%   A20, A21, A24, A26, A27, A35, A40, A41, A51, A57, A58) and a reference to
%   the section that defines it; each Standard ML entry is a fact of the
%   Definition of Standard ML (Revised), read from its contents and body.
% - ILLUMINATING: a reader who knows the Definition of Standard ML can see,
%   construct by construct, where this definition follows it (D7) and where
%   and why it departs from it; no row argues for a design.

\chapter{Elm and Standard ML}\label{ch:appendix-standard-ml}

This appendix compares Standard Elm, as this document defines it, with
Standard ML as defined in \emph{The Definition of Standard ML (Revised)} by
Milner, Tofte, Harper and MacQueen (MIT Press, 1997). It defines nothing; each
entry refers to the place in this document where the construct is defined.

\section{What Is the Same}

The method of this definition is that of the Definition of Standard ML. % D1, D7
The language is presented as a grid of phases (syntax, static semantics,
dynamic semantics) against levels; each semantics section defines its
semantic objects before its rules; the rules infer sentences of the form
background, phrase, meaning, one rule group per phrase class, numbered through
the whole book and followed by Comments; and the rules are given for a bare
language, into which every derived form is translated. Much of the Core
coincides as well: type names generated by type declarations, type functions,
type schemes and their instances, closure of a type with respect to a
context, contexts and environments, function closures whose recursion is
expressed by the operation $\Rec$, matches whose failure is the object
$\FAIL$, value constructors with their identifier status, and evaluation by
value. The differences lie in the effects the dynamic semantics must account
for, in the types of records, tuples and overloaded operations, in the
modules, and in the syntax, and are listed below. % D7

\section{Differences by Construct}

The differences are shown construct by construct in
Figures~\ref{fig:appendix-standard-ml-evaluation} (evaluation),
\ref{fig:appendix-standard-ml-types} and~\ref{fig:appendix-standard-ml-polymorphism}
(types), \ref{fig:appendix-standard-ml-modules} (modules, programs and
effects) and~\ref{fig:appendix-standard-ml-syntax} (syntax). In the figures,
``Ch.'', ``Sec.'', ``Fig.'' and ``App.'' refer to chapters, sections, figures
and appendices of this document.

\newcommand{\smlcmphead}{%
  \textbf{Feature} & \textbf{Standard ML} & \textbf{Elm (this definition)} &
  \textbf{Where defined here}\\ \hline}
\newenvironment{smlcmp}{%
  \small
  \begin{tabular}{@{}>{\raggedright\arraybackslash}p{0.15\textwidth}%
                    >{\raggedright\arraybackslash}p{0.26\textwidth}%
                    >{\raggedright\arraybackslash}p{0.30\textwidth}%
                    >{\raggedright\arraybackslash}p{0.17\textwidth}@{}}
  \smlcmphead}%
  {\end{tabular}}

\begin{figure}[tbp]
\centering
\begin{smlcmp}
Exceptions &
\kw{exception} declarations, \kw{raise} and \kw{handle}; exception names and
packets; an exception convention for the dynamic rules. &
None: no exception names, no packets, no exception convention. Evaluation of
a phrase either gives a value or gives none; absence and failure are values
of the types \kw{Maybe} and \kw{Result} of the initial basis. % D2
& Ch.~\ref{ch:dynamic-core}, Sec.~\ref{sec:dynamic-core-rules};
Sec.~\ref{sec:appendix-static-basis-data}\\ \hline
Incomplete matches &
Allowed; a match that fails raises \kw{Match}, a value binding that fails
raises \kw{Bind}. &
Not allowed: every match, lambda pattern and destructuring binding must be
exhaustive, and every match rule must be needed. &
Sec.~\ref{sec:static-core-further}\\ \hline
State &
References (\kw{ref}, \kw{:=}, \kw{!}), addresses and states; a state
convention, under which the order of premises is the order of evaluation. &
None: sentences carry no state, and no state convention is needed. The order
of premises has no meaning; the order of evaluation is observable only through
\kw{Debug}, and is left to the implementation. % A57, A58
& Sec.~\ref{sec:dynamic-core-rules};
Sec.~\ref{sec:appendix-implementation-order}\\ \hline
Arithmetic failure &
\kw{Overflow} is raised when a result leaves the range of \kw{int}; \kw{Div}
when a divisor is zero. &
Total: every Int result is normalised into the implementation's range, and
an operation with no correct result gives a defined one; thus
\kw{3 // 0} is \kw{0}. % D3, A40, A41
& Ch.~\ref{ch:appendix-dynamic-basis};
Sec.~\ref{sec:appendix-implementation-int}\\ \hline
Equality &
Equality types: type variables \kw{'{}'a}, an equality attribute on type
names; \kw{=} at a function type does not elaborate. &
No equality types: \kw{==} elaborates at every type. When it reaches two
functions its result is an unspecified Bool, so evaluation may be
non-deterministic. % A35, A51
& Sec.~\ref{sec:static-core-classes};
Sec.~\ref{sec:appendix-dynamic-basis-apply};
Sec.~\ref{sec:appendix-implementation-equality}\\
\end{smlcmp}
\caption{Standard ML and Elm: Evaluation}
\label{fig:appendix-standard-ml-evaluation}
\end{figure}

\begin{figure}[tbp]
\centering
\begin{smlcmp}
Overloading &
A side mechanism (Appendix~E of that definition): overloaded identifiers such
as \kw{+} and \kw{<}, and overloaded special constants, resolved by context,
with \kw{int} as default. &
Part of the type system: type variables constrained to the classes
$\knum$, $\kcomp$, $\kapp$ and $\kcapp$, generalised and instantiated
like others. A $\knum$ variable left undetermined in a program is taken to be
Int. % A27
& Sec.~\ref{sec:static-core-classes};
Sec.~\ref{sec:static-core-further}\\ \hline
Records &
Record types are closed. The selector \kw{\#}\lab{} and the record pattern
wildcard \kw{...} require the record type to be determined by the context.
No record update. &
Record types may be extensible, by a row variable
(\kw{\kwlb\ r | x :\ Int \kwrb}). Field selection, the field accessor
\kw{.}\lab{} and record patterns apply to every record with the fields
named. Record update \kw{\kwlb\ r | x = e \kwrb} keeps the type of the record.
% D5
& Sec.~\ref{sec:static-core-compound}; Sec.~\ref{sec:static-core-rules};
App.~\ref{ch:appendix-derived-forms}\\ \hline
Tuples &
A tuple is a record with labels \kw{1}, \ldots, $n$ ($n \geq 2$); the unit is
the empty record. &
Tuples of two and three components have type names of their own, and are not
records; the unit is a type of its own. % A1, A12
& Fig.~\ref{fig:core-grammar-expressions};
Sec.~\ref{sec:static-core-compound}\\ \hline
Datatypes &
\kw{datatype} with \kw{withtype}, datatype replication, \kw{abstype};
datatypes may be declared locally. &
Custom types are declared only at the top level of a module. All custom types
of a module form one group, and they and the type aliases of the module may
mention one another; there is no \kw{withtype}, replication or \kw{abstype}.
A record alias also declares a constructor. &
Fig.~\ref{fig:grammar-top-level-declarations};
Sec.~\ref{sec:appendix-derived-forms-modules};
Sec.~\ref{sec:static-modules-rules}\\
\end{smlcmp}
\caption{Standard ML and Elm: Types}
\label{fig:appendix-standard-ml-types}
\end{figure}

\begin{figure}[tbp]
\centering
\begin{smlcmp}
Value polymorphism &
A value restriction: only a non-expansive expression has its type
generalised. &
No value restriction: closure applies to every binding, whatever its
expression. & Sec.~\ref{sec:static-core-closure};
Sec.~\ref{sec:static-core-rules}\\ \hline
Explicit type variables &
Bound at a value declaration, explicitly (\kw{val 'a}) or implicitly at the
outermost value declaration in which they occur. Type constraints \expm\ \kw{:}\ \ty{}
on any expression or pattern. &
No binder: the type variables of an annotation scope over the annotations
nested in its definition, and may not escape it. Annotations are given only
for definitions and ports. % A20, A21
& Sec.~\ref{sec:static-core-explicit}\\
\end{smlcmp}
\caption{Standard ML and Elm: Polymorphism and Type Variables}
\label{fig:appendix-standard-ml-polymorphism}
\end{figure}

\begin{figure}[tbp]
\centering
\begin{smlcmp}
Modules &
Structures, signatures and functors; sharing; transparent and opaque
signature matching, with realisation and enrichment. &
Flat modules: a module is a header, imports and top-level declarations, and
exports top-level names by an exposing list; a custom type is exposed with or
without its constructors. No signatures, functors or sharing; the import
relation between modules has no cycle. % D6
& Ch.~\ref{ch:syntax-modules}; Ch.~\ref{ch:static-modules},
Sec.~\ref{sec:static-modules-interfaces}; Ch.~\ref{ch:dynamic-modules}\\
\hline
Programs &
An interactive top level: a sequence of top-level declarations, each
elaborated and evaluated in the basis left by those before it; a declaration
that fails is rejected and the session continues. &
A program is a main module and the modules it imports, elaborated as a whole;
its value \kw{main} is executed by a platform, as a labelled transition system
whose observable behaviour is its traffic on ports. &
Ch.~\ref{ch:programs}, Sec.~\ref{sec:programs-execution}\\ \hline
Effects &
At the level of the language, only state and exceptions; input and output are
left to library functions with side effects. &
Evaluation has no effect. Tasks, commands and subscriptions are values, built
by functions of the basis, which the platform carries out; their results
return to the program as messages. % D10
& Ch.~\ref{ch:effects}, Sec.~\ref{sec:effects-tasks},
Sec.~\ref{sec:effects-commands}\\ \hline
Ports &
None. &
A \kw{port} declaration in a port module gives a command for sending values
out or a subscription for receiving them. Only values of port types may cross
a port, encoded as JSON values in which byte sequences may stand. % D11, A92
& Sec.~\ref{sec:static-modules-ports}; Sec.~\ref{sec:effects-ports}\\
\end{smlcmp}
\caption{Standard ML and Elm: Modules, Programs and Effects}
\label{fig:appendix-standard-ml-modules}
\end{figure}

\begin{figure}[tbp]
\centering
\begin{smlcmp}
Infixed operators &
\kw{infix}, \kw{infixr} and \kw{nonfix} declarations in programs, scoped like
other declarations; \kw{op} removes infix status. &
The fixities are fixed by the initial basis and by library modules; no
program declares an operator. Some operators are non-associative. The
operator value \kw{(}\opv\kw{)} uses an operator as a function. % A11
& Sec.~\ref{sec:syntax-core-infix}, Fig.~\ref{fig:core-fixities}\\ \hline
Layout &
Insensitive to layout; phrases are delimited by reserved words such as
\kw{end} and by \kw{;}. &
Sensitive to indentation: match rules, the definitions of a \kw{let} and the
declarations of a module are separated by layout, and \kw{;} and \kw{and} are
abstract syntax. % A5, A6, A17
& Sec.~\ref{sec:syntax-core-layout}; Sec.~\ref{sec:syntax-modules-layout}\\
\hline
Recursion &
\kw{val rec} and \kw{fun} with \kw{and} group recursive bindings as written;
each declaration is in scope only after it. &
The definitions of a \kw{let} and the declarations of a module stand in any
order and are grouped, by derived forms, into components of the ``refers to''
relation in dependency order. In a \kw{let}, recursion passes only through
definitions with parameters. % A24, A26
& Sec.~\ref{sec:appendix-derived-forms-let};
Sec.~\ref{sec:appendix-derived-forms-modules};
Sec.~\ref{sec:syntax-core-restrictions}\\ \hline
Functions &
\kw{fn} takes a match; \kw{fun} defines a function by several clauses. &
A lambda expression takes atomic patterns and has one body; a function
binding has one clause; several cases are written with \kw{case}. &
Fig.~\ref{fig:core-grammar-expressions};
App.~\ref{ch:appendix-derived-forms}\\
\end{smlcmp}
\caption{Standard ML and Elm: Syntax}
\label{fig:appendix-standard-ml-syntax}
\end{figure}

Standard ML constructs without a counterpart here are the sequential
expression, \kw{while}, \kw{local}, \kw{open}, \kw{abstype} and the
declarations of exceptions; the nearest to \kw{open} is the exposing list of
an import (Section~\ref{sec:syntax-modules-grammar}). Elm constructs without a
counterpart in Standard ML are record update and extensible record types,
the field accessor, the constrained classes of type variables, and ports.
